% Generated by build_tex_sections.py; full proof text, no abridgement.
% Requires amsmath, amssymb, and UTF-8 support.
\section{A rank-zero semiregular complex at a CM sixfold}

Run HODGE\_DIVISION; agent \texttt{/root/imprimitive}; 2026-11-21. Model GPT-6 Astra Pro, effort inherited/not exposed. \textbf{D-SR-01--06 passed independent whole-proof checks by \texttt{/root/sixfolds}; D-SR-01--05 also passed a separate whole-proof audit by \texttt{/root/g34\_eightfolds}.} See the peer reports and \texttt{CHECK\_EVENTS.tsv} for the exact scope of each event. No result below relies on a theorem of Markman. His secant construction suggested the approach. The classical trace identities, Fourier--Mukai equivalence and its trace naturality are separate cited inputs. The final deformation implication is exactly the supplied DEG\_L S-14/Perry theorem; its preprint status is retained. The later positive-rank construction SF-QC01--05 supplies an explicit numerical detector bound; the present section retains the independent rank-zero construction and its integral normalization. \cite{SR-Perry2026}

\subsection{D-SR-01. A finite construction of a disjoint cyclic orbit}
\label{D-SR-01}

Let \(C\) be a smooth nonhyperelliptic genus-three curve, \(J=J(C)\), and \(\Theta\) its principal polarization. Put \(D=C-C\) in \(J\). Then \(D\) is an irreducible divisor of class \(2\Theta\). Indeed the difference map \(C\times C\) to \(D\) is generically one-to-one: a nontrivial equality \(p-q=p^{\prime}-q^{\prime}\) would yield two different effective degree-two divisors in the same linear system, contradicting nonhyperellipticity. The Poincare formula and the Pontryagin product give \([C]*[-C]=2\Theta\). We use the usual Abel--Jacobi embedding; translating it does not alter \(D\). \cite{SR-BirkenhakeLange}

For every positive integer \(r\),

\[
\#\bigl(D\cap J[r]\bigr)\leq 108r^4.\tag{SR1.1}
\]

Here points are counted without multiplicity. To prove this, embed \(J\) by \(3\Theta\). Hyperplanes through the image of zero have base locus exactly \{0\}. Their pullbacks by [\(r\)] have finite base locus \(J[r]\). Choose two general such hyperplanes so that their pullbacks cut the surface \(D\) properly. Each point of \(D\) intersect \(J[r]\) contributes at least one to the resulting zero-dimensional intersection. Its degree is \cite{SR-BauerSzemberg}

\[
D\cdot([r]^*3\Theta)^2=(2\Theta)\cdot(3r^2\Theta)^2=18r^4\Theta^3=108r^4.
\]

The existence of proper successive cuts follows from Bertini avoidance of the finitely many positive-dimensional components; smoothness of the cuts is unnecessary.

Set

\[
P=2^{2048},\quad L=\operatorname{lcm}(1,2,\ldots,P),\quad N=4L^2+1,\quad d=N-1=4L^2.\tag{SR1.2}
\]

There is a cyclic subgroup \(H\) of \(J[N]\), of exact order \(N\), such that

\[
H\cap(C-C)=\{0\}.\tag{SR1.3}
\]

Here is a finite counting proof, with an explicit search specification. Choose \(v\) uniformly among the points of exact order \(N\) in \(J[N]=(\mathbb Z/N)^6\). For every divisor \(q> 1\) of \(N\), the point \((N/q)v\) is uniformly distributed among the exact-order-q points. Their number is

\[
J_6(q)=q^6\prod_{p\mid q}(1-p^{-6})>\frac{q^6}{2}.
\]

The inequality follows, for example, from \(\prod_i(1-a_i)\geq1-\sum_i a_i\) and \(\sum_{n\geq2}n^{-6}<1/2\). If the cyclic subgroup generated by \(v\) has a nonzero point on \(D\), one of its exact-order-q generators lies on \(D\) for some \(q\mid N\). The union bound and (SR1.1) therefore give

\[
\begin{aligned}\Pr\bigl(\langle v\rangle\cap(D\setminus\{0\})\ne\varnothing\bigr)&\leq216\sum_{\substack{q\mid N\\q>1}}\frac{\varphi(q)}{q^2}\\&\leq216\left[\exp\left(\sum_{p\mid N}\frac1p\right)-1\right].\end{aligned}\tag{SR1.4}
\]

The first estimate allows all \(\varphi(q)\) generators of the cyclic order-q subgroup and thus overcounts, harmlessly. The second uses

\[
\sum_{q\mid N}\frac{\varphi(q)}{q^2}=\prod_{p^a\parallel N}\left(1+\frac{1-p^{-a}}p\right)\leq\prod_{p\mid N}\left(1+\frac1p\right).
\]

Every prime divisor of \(N\) exceeds \(P\), because every prime \textless{}=\(P\) divides \(L\) and \(N\) is 1 modulo that prime. For \(P\) a power of two, the elementary central-binomial argument gives

\[
\log L<P\log4.
\]

For completeness, if \(\psi(x)=\log\operatorname{lcm}(1,\ldots,\lfloor x\rfloor)\), then \(\psi(2r)-\psi(r)\leq\log\binom{2r}{r}\leq2r\log2\): every prime power in (\(r\),2r] contributes its prime once to the binomial coefficient. Sum over the dyadic intervals ending at \(P\). Consequently

\[
\sum_{p\mid N}\frac1p\leq\frac{\log N}{(P+1)\log(P+1)}<\frac1{500}.
\]

The last bound follows from \(\log N\leq\log5+2\log L<\log5+4P\log2\) and \(\log P=2048\log2\). Thus (SR1.4) is strictly less than 216/499\textless{}1/2, using \(\exp(x)-1\leq x/(1-x)\) for 0\textless{}\(x\)\textless{}1. This proves existence.

To specify a finite construction, fix any algebraic enumeration of the finite group \(J[N]\), test its exact-order-N points in that order, and take the first \(v\) for which kv is outside the closed algebraic divisor \(D\) for every 1\textless{}=k\textless{}\(N\). All these tests are finite algebraic membership tests, and the proof guarantees termination. This is an explicit terminating construction, although its deliberately coarse integer \(N\) is extremely large. It is not a claim that we have enumerated \(J[N]\) in this run.

Let \(q:J\to Z=J/H\) be the quotient. By (SR1.3), the restriction \(q|_C\) is a closed embedding: it is injective and unramified, and the source is proper. Write \(\bar C=q(C)\). Its full inverse image is the disjoint union of the \(N\) translates \(C\)+\(h\).

Let \(\theta\) in \(\operatorname{NS}(Z)_{\mathbb Q}\) be the unique rational class with \(q^*\theta=\Theta\). It is ample. Uniqueness and existence follow from the rational cohomology isomorphism of an isogeny and the norm of a line bundle; explicitly \(q^*q_*\Theta=N\Theta\) because translations act trivially on cohomology. Put \(F=\mathcal I_{\bar C}\). Then

\[
\operatorname{ch}(F)=1-\frac N2\theta^2+\frac N3\theta^3,\qquad\Delta:=\int_Z\frac{\theta^3}{6}=\frac1N.\tag{SR1.5}
\]

The character follows by pulling back to \(J\): a smooth genus-three curve has \(\operatorname{ch}(\mathcal O_C)=[C]-2[\mathrm{pt}]\), since the degree of its Euler characteristic is -2 and \(\operatorname{td}(J)=1\). Pullback of the ideal is the ideal of the disjoint union, so its character is \(1-N\Theta^2/2+2N[\mathrm{pt}]\).

\textbf{Extension.} For any fixed positive integer \(m\leq 163\) the same proof applies to \(N=4mL^2+1\) and \(d=4mL^2\), with the same \(P\),\(L\). Replace log5 by log(4m+1) in the final estimate; it still gives \(\sum_{p\mid N}1/p<1/500\). The resulting field below is \(\mathbb Q(\sqrt{-m})\). This extension concerns a CM sixfold isogenous to the Fermat Jacobian square; for \(m\ne 1\) it is not being identified with a power of an elliptic curve with CM by \(\mathbb Q(\sqrt{-m})\).

\subsection{D-SR-02. The descended ideal and its external square are semiregular}
\label{D-SR-02}

In the situation above, \(F=\mathcal I_{\bar C}\) on the abelian threefold \(Z\) has

\[
\bigl(\dim\operatorname{Ext}^j(F,F)\bigr)_{j=0}^3=(1,6,6,1).\tag{SR2.1}
\]

Its trace maps \(\sigma_F^j:\operatorname{Ext}^j(F,F)\to\operatorname{HH}_{-j}(Z)\) are injective for \(j= 0\),1,2. The same statements hold for the derived dual \(F^\vee\). Consequently \(F\boxtimes F^\vee\) on \(Z\times Z\) is a gluable fully semiregular perfect complex, with \(\operatorname{Ext}^2\) dimension 48.

\textbf{Proof of (SR2.1).} Since \(q\) is etale and \(q|_C\) is an isomorphism onto \(\bar C\), the normal bundles \(N_{C/J}\) and \(N_{\bar C/Z}\) are identified. The normal sequence

\[
0\longrightarrow T_C\longrightarrow\mathcal O_C^3\longrightarrow N_{C/J}\longrightarrow0.
\]

shows \(h^0(N)=3\). Indeed the next map \(H^1(T_C)\to H^1(\mathcal O_C)^3\) is injective: its dual is the multiplication \(H^0(K_C)\otimes H^0(K_C)\to H^0(K_C^2)\), which is surjective for a smooth plane quartic. In this genus-three case this last statement is elementary: \(K_C=\mathcal O_C(1)\), and restriction of plane quadrics is an isomorphism from the six-dimensional space of quadrics onto \(H^0(C,\mathcal O_C(2))\); no nonzero quadric contains \(C\). Also deg \(N= 4\) and rank \(N= 2\), so Riemann--Roch gives \(h^1(N)=3\).

Here is a direct self-Ext calculation, avoiding any global identification of moduli spaces. Apply Hom(\(F\),-) to \(0\to F\to\mathcal O_Z\to\mathcal O_{\bar C}\to0\). The map Hom(\(F\),O\_Z)-\textgreater{}\(\operatorname{Hom}(F,\mathcal O_{\bar C})\) is zero, since every global homomorphism \(F\to\mathcal O_Z\) is a scalar multiple of the ideal inclusion. Thus there is an exact sequence

\[
0\longrightarrow H^0(N_{\bar C/Z})\longrightarrow\operatorname{Ext}^1(F,F)\longrightarrow\operatorname{Ext}^1(F,\mathcal O_Z).
\]

By threefold Serre duality, \(\operatorname{Ext}^1(F,\mathcal O_Z)=H^2(F)^*\). The restriction \(H^1(\mathcal O_Z)\to H^1(\mathcal O_{\bar C})\) is an isomorphism: it is the Abel--Jacobi isomorphism after composition with the isogeny \(q\). The cohomology sequence of the ideal therefore gives \(H^2(F)=H^2(\mathcal O_Z)\), of dimension three. Thus dim \(\operatorname{Ext}^1(F,F)\leq3+3=6\). The three embedded translation directions and the three degree-zero tensor directions are independent. Indeed the embedded ideal deformations have trivial determinant, whereas the trace of a tensor direction is the original class in \(H^1(\mathcal O_Z)\), because \(\operatorname{rank}(F)=1\). Hence equality holds and

\[
\operatorname{Ext}^1(F,F)=H^1(\mathcal O_Z)\oplus H^0(N_{\bar C/Z}).\tag{SR2.2}
\]

These assertions can also be checked on the two-generator Koszul resolution of the ideal; no stability or smoothness of a moduli space of sheaves is assumed.

\(\operatorname{End}(F)=\mathbb C\): restriction outside \(\bar C\) identifies an endomorphism with a regular function, which extends across codimension two on the normal projective variety \(Z\). Serre duality, with \(\omega_Z=\mathcal O_Z\), then yields (SR2.1).

Write \(\sigma_F^j(x)=\operatorname{tr}(\exp(\operatorname{at}(F))x)\) under HKR; the Todd correction is one for an abelian variety. For \(j= 0\) the trace of the identity is \(\operatorname{ch}(F)\), whose rank component is 1. For \(j= 1\), the infinitesimal action of \(H^1(\mathcal O_Z)\) by tensoring and of \(H^0(T_Z)\) by translation gives an isomorphism \cite{SR-BuchweitzFlenner2003,SR-BuchweitzFlenner2008}

\[
H^1(\mathcal O_Z)\oplus H^0(T_Z)\longrightarrow\operatorname{Ext}^1(F,F).
\]

The translation summand is an isomorphism onto \(H^0(N)\) by the normal sequence just used, and the tensor summand maps isomorphically to the determinant direction. Composition with \(\sigma_F^1\) sends the first summand to \(H^1(\mathcal O_Z)\) by the identity; the \(H^0(T_Z)\) summand has zero component there because \(c_1(F)=0\) and maps injectively into \(H^2(\Omega_Z^1)\) by contraction with \(-(N/2)\theta^2\). Nondegeneracy of \(\theta\) proves that contraction is injective. Thus \(\sigma_F^1\) is injective.

For \(j= 2\), apply the trace/Atiyah obstruction identity to the subspace

\[
H^2(\mathcal O_Z)\oplus H^1(T_Z)\subset\operatorname{HH}^2(Z).
\]

The \(H^2(\mathcal O_Z)\) summand maps by the rank-one identity to the \(H^2(\mathcal O_Z)\) component of \(\operatorname{HH}_{-2}\). The \(H^1(T_Z)\) summand has zero component there because \(c_1(F)=0\) and maps surjectively onto \(H^3(\Omega_Z^1)\) by contraction with \(-(N/2)\theta^2\). Here is a direct rank check. In holomorphic/antiholomorphic bases \(\theta=\sum_i\alpha_i\beta_i\), contraction by a Kodaira--Spencer matrix \(\kappa\) gives \(\kappa\) contraction \(\theta^2\)=2(\(\kappa\) contraction \(\theta\))\(\theta\). The map kappa-\textgreater{}\(\kappa\) contraction \(\theta\) is onto \(\bigwedge^2\langle\beta_i\rangle\), and multiplication by \(\theta\) is an isomorphism of that three-dimensional space onto \(\langle\alpha_i\rangle\otimes\bigwedge^3\langle\beta_i\rangle\). Hence the rank is three. Therefore the full composite has rank six. Its middle space \(\operatorname{Ext}^2(F,F)\) has dimension six, so \(\sigma_F^2\) is injective. Only the classical trace/Atiyah identity is used, not a conjectural deformation assertion.

Derived duality identifies the self-Ext groups of \(F\) and \(F^\vee\); the same rank argument applies to \(\operatorname{ch}(F^\vee)=1-(N/2)\theta^2-(N/3)\theta^3\). In particular \(\sigma_{F^\vee}^j\) is injective for \(j= 0\),1,2.

Finally the trace of an external tensor product is the external product of the traces, because its Atiyah class is \(\operatorname{at}(F)\) external\_tensor id + id external\_tensor \(\operatorname{at}(F^\vee)\), and the exponential and trace factor. The \(\operatorname{Ext}^2\) Kunneth decomposition is

\[
\begin{aligned}&\operatorname{Ext}^2(F,F)\otimes\operatorname{Ext}^0(F^\vee,F^\vee)\\{}\oplus{}&\operatorname{Ext}^1(F,F)\otimes\operatorname{Ext}^1(F^\vee,F^\vee)\\{}\oplus{}&\operatorname{Ext}^0(F,F)\otimes\operatorname{Ext}^2(F^\vee,F^\vee).\end{aligned}
\]

The dimensions are 6,36,6. Their trace images lie in the three distinct Hochschild Kunneth components \(\operatorname{HH}_{-2}\otimes\operatorname{HH}_0\), \(\operatorname{HH}_{-1}\otimes\operatorname{HH}_{-1}\), and \(\operatorname{HH}_0\otimes\operatorname{HH}_{-2}\). Each restriction is injective by the preceding checks, so there is no cancellation between the three pieces. This proves semiregularity and the dimension 48. Negative self-Ext groups vanish by the same Kunneth decomposition, since they vanish for a coherent sheaf and for its derived dual. This proves gluability.

\subsection{D-SR-03. An exact Fourier--Mukai formula with a nonzero Weil term}
\label{D-SR-03}

Let \(A=Z\times\widehat Z\). Let Poincare be the normalized Poincare bundle and \(\gamma=c_1(\mathcal P)\) on A. Let \(\widehat\theta\) be the rational dual polarization satisfying \(\phi_\theta^*\widehat\theta=\theta\), where \(\phi_\theta:Z\to\widehat Z\) is the rational polarization quasi-isogeny. Define the positive rational divisor class

\[
H=N\theta_x+\widehat\theta+\gamma.\tag{SR3.1}
\]

Let \(\mu:Z\times Z\to Z\times Z\) be \(\mu(x,y)=(x+y,y)\). Define the exact equivalence

\[
\Phi=(\operatorname{id}_Z\boxtimes\operatorname{FM}_{\mathcal P^{-1}})\circ\mu^*,\qquad E=\Phi(F\boxtimes F^\vee).\tag{SR3.2}
\]

No extra cohomological shift is included in this convention. Then \(E\) is gluable and fully semiregular, and its complete Chern character satisfies

\[
\exp(-\gamma/2)\operatorname{ch}(E)=t+2H-\frac{H^3}{12d}+\frac{H^5}{960d^2}.\tag{SR3.3}
\]

where \(t\) is a nonzero rational Weil class for a \(K=\mathbb Q(\sqrt{-d})\) action of signature (3,3) on (A,\(H\)). Thus rank(\(E\))=0 causes no problem for the exact S-14 hypothesis.

\textbf{Exact calculation.} Put \(u=\sqrt{-d}\), \(s\)=-1+\(u\), \(\bar s\)=-1-u, and

\[
a=\frac{1+1/u}{2},\qquad\bar a=\frac{1-1/u}{2}.
\]

Equation (SR1.5) is exactly

\[
\operatorname{ch}(F)=a\exp(s\theta)+\bar a\exp(\bar s\theta).\tag{SR3.4}
\]

It suffices to expand through \(\theta^3\); use \(u^2=-d\) and \(N=d+1\). The coefficient of \(\theta\) is zero, that of \(\theta^2\) is -N/2, and that of \(\theta^3\) is \(N\)/3.

Choose rational symplectic Darboux coordinates for \(\theta\). In one of its three planes denote the first-factor coordinates by \(a_i\),\(b_i\), the integrated second-factor coordinates by \(c_i\),\(d_i\), and the dual output coordinates by \(u_i\),\(v_i\). Then \(\theta_x=\sum_i a_i b_i\), \(\theta_y=\sum_i c_i d_i\), \(\widehat\theta=\sum_i u_i v_i\), and \(\gamma=\sum_i(a_i v_i-b_i u_i)\). Juxtaposition means exterior multiplication. The Poincare kernel on the integrated factors has class \(\sum_i(c_i v_i-d_i u_i)\). Direct extraction of the \(c_i\) \(d_i\) coefficient gives

\[
\begin{aligned}&\int_{c_i,d_i}\exp\bigl(s(a_i+c_i)(b_i+d_i)-t c_i d_i-(c_i v_i-d_i u_i)\bigr)\\&\quad=(s-t)-st a_i b_i-s b_i u_i-u_i v_i+s a_i v_i-s a_i b_i u_i v_i.\end{aligned}\tag{SR3.5}
\]

The integration over all three planes has the additional scalar \(\Delta\)=integral\_Z \(\theta^3/6=1/N\). This scalar is essential because \(\theta\) is rational and need not be principal on \(Z\).

When \(t\)=\(s\), (SR3.5) equals

\[
-\exp\bigl((a_i v_i-b_i u_i)/2\bigr)(s a_i-u_i)(s b_i-v_i).
\]

Consequently, putting \(Q_s=\prod_{i=1}^3(s a_i-u_i)(s b_i-v_i)\), the two equal-root terms in the normalized transform are

\[
t=-\Delta\bigl(a^2Q_s+\bar a^2Q_{\bar s}\bigr).\tag{SR3.6}
\]

When \(t\)=\(\bar s\), set \(\alpha=s-\bar s=2u\). Since \(st=N\) and \(s+t=-2\), (SR3.5) is

\[
\alpha\exp(\gamma_i/2-H_i/\alpha),\qquad H_i=N a_i b_i+u_i v_i+a_i v_i-b_i u_i.
\]

The unequal-root pair therefore contributes after normalization

\[
\Delta a\bar a\alpha^3\bigl[\exp(-H/\alpha)-\exp(H/\alpha)\bigr].
\]

Use \(a\bar a=N/(4d)\), \(\alpha^2=-4d\) and \(H^7=0\). Expanding the odd powers gives precisely the polynomial in (SR3.3). Equations (SR3.4)--(SR3.5), or a sparse exterior-algebra expansion of all twelve output generators, are exact identities; a separate script checks the displayed formulas and records the coefficients.

\textbf{The Weil structure and positivity.} The rational endomorphism \(T\) of A is defined, after the quasi-isogeny identification A\textasciitilde{}\(Z\times Z\) supplied by \(\theta\), by the matrix

\[
T=\begin{pmatrix}-2&-1\\N&0\end{pmatrix}.
\]

It satisfies \(T^2+2T+N=0\). Thus \(T+1\) defines \(\sqrt{-d}\). On \(H^1\), its s-eigenspace is exactly the six-dimensional span of the forms \(s\) a\_i-u\_i, \(s\) b\_i-v\_i. Its conjugate eigenspace is the analogous \(\bar s\) span. Since the two-by-two matrix acts by scalars on each copy of \(H^1(Z)\), each eigenspace has holomorphic dimension three and antiholomorphic dimension three. Thus \(Q_s\) and \(Q_{\bar s}\) span the two determinant lines of the rational Weil plane.

The class \(H\) is ample. Under id times \(\phi_\theta\) its pullback is

\[
N\theta_1+\theta_2+(m^*\theta-\theta_1-\theta_2)=(N-1)\theta_1+m^*\theta,
\]

which is ample because (\(x\),\(y\))-\textgreater{}(\(x\),\(x\)+\(y\)) is an isomorphism and \(N> 1\). The generator \(T+1\) is Rosati anti-self-adjoint for \(H\): in these coordinates the polarization matrix is \(S=\bigl(\begin{smallmatrix}N&1\\1&1\end{smallmatrix}\bigr)\), and \(S(T+1)=\bigl(\begin{smallmatrix}0&-d\\d&0\end{smallmatrix}\bigr)\) is skew-symmetric. Hence the \(K\) action is compatible with \(H\). The coefficients a and \(\bar a\) are nonzero, and \(Q_s\),\(Q_{\bar s}\) are independent nonzero determinant vectors, so (SR3.6) is nonzero. Galois conjugation interchanges its two summands, proving rationality. This proves the entire character identity and its nonzero Weil component.

Fourier--Mukai equivalence preserves negative self-Ext and injectivity of the categorical trace, so D-SR-02 proves the stated gluability and semiregularity of \(E\). This preservation is the functoriality in Perry Remarks 2.3 and 2.5, and follows directly from transporting the evaluation/coevaluation defining the categorical trace across an equivalence. The result is not inferred from equality of K-classes alone. \cite{SR-Perry2026}

\subsection{D-SR-04. Transport to the requested balanced CM power}
\label{D-SR-04}

Choose \(C\) to be the Fermat quartic \(x^4+y^4+z^4=0\). Its Jacobian is isogenous to \(E_i^3\), where \(E_i=\mathbb C/(\mathbb Z+i\mathbb Z)\): Beauville, *Some surfaces with maximal Picard number*, JEP 1 (2014), section 4, Example 3, \(p\).107, with Lemma 2 on \(p\).106. Thus \(A=Z\times\widehat Z\) is isogenous to \(E_i^6\). \cite{SR-Beauville2014}

The isogeny to the CM power and its polarization can be specified explicitly. The three degree-two quotients of the Fermat quartic by sign changes of its coordinates are elliptic curves with j=1728. Choose origin-preserving isomorphisms of them with \(E_i\), and let

\[
j:E_i^3\longrightarrow J.
\]

be the sum of the three pullback homomorphisms of Jacobians. Their cotangent images are the three distinct nontrivial characters of the sign-change Klein four group, so j is an isogeny. Each diagonal norm-pullback composition is multiplication by two; cross compositions vanish, since they intertwine different sign characters. Therefore

\[
j^*\Theta=2\theta_{E_i^3},\qquad\deg j=8.
\]

Put \(h=qj:E_i^3\to Z\), \(r=2L\), \(d=r^2\), \(s=-1+ri\), and define the two-by-two matrices acting on threefold pairs by

\[
R=\begin{pmatrix}1&1\\\bar s&s\end{pmatrix},\quad k=\operatorname{diag}(h,\phi_\theta h),\quad f=NkR:E_i^6\longrightarrow A.\tag{SR4.1}
\]

The map k is initially a rational homomorphism, but \(f\) is an actual isogeny: Ntheta is the integral norm-polarization class \(q_*\Theta\), so \(N\phi_\theta=\phi_{N\theta}\) is an actual homomorphism. All entries of \(R\) are Gaussian integers. Its two columns are eigenvectors of \(T+1\) with eigenvalues ri and -ri. Thus \(f\) is K-linear for the specified balanced source action \(\operatorname{diag}(i,i,i,-i,-i,-i)\). If \(\theta_{\rm prod}\) is the product principal polarization on \(E_i^6\), direct matrix multiplication gives

\[
R^\dagger\begin{pmatrix}N&1\\1&1\end{pmatrix}R=2dI,\qquad f^*H=4dN^2\theta_{\rm prod},\qquad\deg f=4096d^3N^{12}.\tag{SR4.2}
\]

For the degree calculation, \(\deg h=8N\), \(\deg\phi_\theta=\Delta^2=N^{-2}\) as a rational isogeny, so \(\deg k=64\); \(\deg R=\operatorname{Norm}(2ri)^3=(4d)^3\); multiplication by \(N\) on a sixfold has degree \(N^{12}\). Multiplying gives the stated integer.

Let \(G=\ker f\), acting honestly by translations, and \(E_0=f^*E\) with its canonical coherent G-linearization. Etale descent identifies

\[
\operatorname{Ext}^2(E_0,E_0)^G=\operatorname{Ext}^2(E,E).
\]

Under HKR, the full trace of a pullback along an etale morphism is the pullback of the trace. The map \(f^*:\operatorname{HH}_{-2}(A)\to\operatorname{HH}_{-2}(E_i^6)\) is injective; indeed each Hodge-cohomology pullback is an isomorphism over Q or \(C\) for an isogeny, with inverse \(f_*/\deg f\). Since \(E\) is fully semiregular, the trace of \(E_0\) is injective on the invariant \(\operatorname{Ext}^2\). Hence \(E_0\) is G-semiregular, and therefore weakly G-semiregular.

We check ordinary negative self-Ext, rather than only its G-invariant part. The finite etale isogeny has

\[
f_*\mathcal O_{E_i^6}=\bigoplus_{\alpha\in\ker\widehat f}\alpha,
\]

where the summands are torsion degree-zero line bundles on A. Adjunction and the projection formula yield

\[
\operatorname{Ext}^j(E_0,E_0)=\bigoplus_\alpha\operatorname{Ext}^j(E,E\otimes\alpha).
\]

Conjugating tensorization by \(\alpha\) through \(\Phi\) gives an autoequivalence of \(Z\times Z\) consisting of a translation followed by tensorization with a degree-zero line bundle. This can be seen directly from (SR3.2): Fourier--Mukai interchanges a degree-zero tensor on the dual factor with a translation on the original factor, and \(\mu\) conjugates translations to translations and pulls degree-zero line bundles back to degree-zero line bundles. Every such operation on \(Z\times Z\) is a product operation, since \(\operatorname{Pic}^0(Z\times Z)=\operatorname{Pic}^0(Z)\times\operatorname{Pic}^0(Z)\). Thus, up to an irrelevant one-dimensional scalar factor, \cite{SR-Mukai1981}

\[
\Phi^{-1}(E\otimes\alpha)=(t_x^*F\otimes L_1)\boxtimes(t_y^*F^\vee\otimes L_2).
\]

The self-Ext comparison with \(F\boxtimes F^\vee\) has no negative degrees: the first factor is Ext between coherent sheaves; by derived duality the second is also Ext between coherent sheaves in the opposite order. Their Kunneth tensor product therefore has no negative cohomology. This proves \(\operatorname{Ext}^j(E_0,E_0)=0\) for every \(j< 0\), closing the gluability requirement.

The character identity already pulls back exactly: set \(B_0=-f^*\gamma/2\), \(t_0=f^*t\), and \(H_0=f^*H\). Then

\[
\exp(B_0)\operatorname{ch}(E_0)=t_0+2H_0-\frac{H_0^3}{12d}+\frac{H_0^5}{960d^2}.
\]

By (SR4.2), \(H_0=4dN^2\theta_{\rm prod}\). We can take the fixed very ample polarization \(\theta_0=3\theta_{\rm prod}\) and rewrite the right side as \(t_0+P(\theta_0)\), with the explicitly displayed rational coefficients after substituting \(H_0=(4dN^2/3)\theta_0\). Thus this is a seed on the requested balanced CM power with its product principal polarization, not merely on an unspecified isogenous polarized variety. All of the S-14 object hypotheses have now been proved and independently checked.

\subsection{D-SR-05. Split determinant and the other imaginary quadratic fields}
\label{D-SR-05}

The K-Hermitian space in D-SR-03 is split, of rank six and Witt index three; in the raw determinant convention its determinant class is \([-1]\). This does not follow merely from its real signature. Here is a direct rational proof. Let \(E_\theta\) be the rational alternating polarization form on \(V=H_1(Z,\mathbb Q)\), and write \(J\)=\(T+1\) and \(u=\sqrt{-d}\). Define \(q_H(x,y)=E_H(x,Jy)\) and \(h(x,y)=q_H(x,y)+uE_H(x,y)\). Since \(J\) is skew-adjoint and \(J^2=-d\), \(q_H\) is symmetric, \(h\) is Hermitian, \(h(x,Jy)=u h(x,y)\), and \(h(Jx,y)=-u h(x,y)\). Thus \(h\) is conjugate-linear in the first argument and linear in the second, with \(h(v,v)=q_H(v,v)\).

Choose any rational E\_theta-Lagrangian \(L\) in \(V\), of dimension three. Under the rational identification \(H_1(A,\mathbb Q)=V\oplus V\), put \(I=L\oplus L\). The scalar two-by-two matrix \(J\) preserves I, so I has K-dimension three. The polarization matrix is \(S\otimes E_\theta\); hence \(E_H\) vanishes on I. Because I is J-stable, \(q_H\) also vanishes on I, and therefore \(h\) vanishes on I. This is a half-dimensional totally isotropic K-subspace. A nondegenerate Hermitian space with such a subspace is a sum of three hyperbolic planes: choose a dual basis for I, and subtract half of its Hermitian self-pairing matrix to make the dual subspace isotropic. Each hyperbolic plane has determinant -1 modulo K-norms. Thus \(h\) is split and has determinant \((-1)^3\) modulo K-norms. This argument concerns the actual polarization and \(K\) action and works independently of the additional CM endomorphisms.

For any positive squarefree integer \(m\), choose

\[
\begin{gathered}P_m=2^{2048+\lceil\log_2(m+1)\rceil},\qquad L_m=\operatorname{lcm}(1,\ldots,P_m),\\N_m=4mL_m^2+1,\qquad d_m=4mL_m^2.\end{gathered}
\]

The proof of D-SR-01 applies verbatim. Indeed every prime of \(N_m\) exceeds \(P_m\), log \(L_m<P_m\) log4, and

\[
\sum_{p\mid N_m}\frac1p<\frac{\log(4m+1)+4P_m\log2}{P_m\log P_m}<\frac1{500}.
\]

For the last inequality use \(m<P_m\), log \(P_m\geq2048\log2\) and \(P_m\geq2^{2048}\); the two lower-order logarithmic terms are less than \(2/P_m\). Thus the same strict union bound gives the required cyclic subgroup. The construction D-SR-01--03 yields a gluable fully semiregular complex on a CM sixfold isogenous to \(E_i^6\), with the complete character identity and Weil field \(K\)=\(\mathbb Q(\sqrt{-m})\). For the nine class-number-one fields the smaller common \(P=2^{2048}\) already suffices, as stated in D-SR-01. The underlying elliptic CM field is \(\mathbb Q(i)\); the newly constructed Weil action is \(K\). Those two actions must not be confused.

Apply supplied DEG\_L S-14 with \(G\) trivial on A and with a sufficiently divisible integral multiple of \(H\) as the relatively very ample polarization. On a finite level cover of the corresponding split Weil family, the rational Weil plane is a constant local system; take the flat section whose value at this CM point is \(t\). Formula (SR3.3) is its complete Hodge-compatible character. S-14 therefore gives bounded division for this nonzero Weil class throughout that fixed family. The supplied Theorem W then gives the corresponding detector/whole-Weil-plane conclusion. This statement uses Perry's deformation theorem as specified in S-14; it has no Markman, Severn, or 032 theorem as a premise. \cite{SR-Perry2026}

The construction settles the existence of the deformation object, including at the standard balanced Gaussian CM power in D-SR-04. It supplies an explicit finite search bound for the cyclic subgroup, an exact Chern character, an actual group action, and an exact CM isogeny degree. D-SR-06 additionally supplies an ordinary, untwisted deformation and the explicit class denominator and norm. This rank-zero argument by itself leaves the detector degree existential. The later positive-rank construction SF-QC01--05 supplies the missing effective presentation and a numerical degree bound; it has a separate full proof and audit. A positive definite form or a divisor-generated CM K-class alone does not supply that presentation.

\subsection{D-SR-06. An integral twist, ordinary deformation, and exact class normalization}
\label{D-SR-06}

Retain the Gaussian CM isogeny \(f\) of D-SR-04 and put \(c=4dN^2\), so \(H_0=c\theta_{\rm prod}\). The class \(B_0\) is integral. To check this with the correct normalization, let \(C_2=\bigl(\begin{smallmatrix}0&1\\1&0\end{smallmatrix}\bigr)\). The pullback of \(\gamma\) by \(k=\operatorname{diag}(h,\phi_\theta h)\) has matrix \(2C_2\) relative to the product principal polarization, because \(h^*\theta=2\theta_{E_i^3}\). Consequently

\[
\begin{gathered}R^\dagger C_2R=2\begin{pmatrix}-1&s\\\bar s&-1\end{pmatrix},\quad f^*\gamma=4N^2\begin{pmatrix}-1&s\\\bar s&-1\end{pmatrix}\otimes I_3,\\B_0=2N^2\begin{pmatrix}1&-s\\-\bar s&1\end{pmatrix}\otimes I_3.\end{gathered}\tag{SR6.1}
\]

The entries of the last Hermitian matrix are Gaussian integers and its diagonal entries are ordinary integers. Such a matrix represents an integral divisor class on the principally polarized Gaussian product. Indeed the map from divisor classes to symmetric endomorphisms, composed with the inverse principal polarization, identifies these classes with the integral Hermitian matrices; equivalently, its alternating Riemann form is integral on the product Gaussian lattice. By the integral (1,1) theorem it is the first Chern class of a line bundle \(L_B\). No ampleness of \(L_B\) is needed.

Put \(E_0^{\prime}=E_0\otimes L_B\). Conjugate the honest G-action by tensorization with \(L_B\). For a translation \(t_g\) this replaces the action by \(t_g^*(-)\) tensor (\(L_B\) tensor \(t_g^*L_B^{-1}\)). The additional line bundle is degree zero, and the assignment is a homomorphism into the translation/tensor group; its image is finite because \(G\) is finite. The coherent composition isomorphisms are transported by the tensor equivalence. The equivariant structure, weak semiregularity, and ordinary negative-self-Ext vanishing are therefore preserved. Now the entire ordinary character is

\[
\operatorname{ch}(E_0^{\prime})=t_0+2c\theta_{\rm prod}-\frac{16}{3}d^2N^6\theta_{\rm prod}^3+\frac{16}{15}d^3N^{10}\theta_{\rm prod}^5.\tag{SR6.2}
\]

All three polynomial coefficients are integers: \(L\) is divisible by 3 and 5, and \(d=4L^2\). This is a genuine line-bundle twist of the complex; it is not a formal normalization of a rank-zero character.

In Perry's proof of Theorem 1.2 in section 6.2, pp.26--28, one can now choose \(B_0^{\prime}=0\), hence the initial Kummer class and the associated Brauer class are zero. The band pushout used to extend the group action sends the zero gerbe class to zero. The ambient category in Theorem 6.3 is consequently Perf of the ordinary abelian family. Thus the resulting object on an etale neighborhood is an ordinary perfect complex, even if a separate presentation of the invariant category involves a gerbe. This conclusion uses the actual proof of the cited preprint, rather than strengthening the word "twisted" in its introductory statement without justification. \cite{SR-Perry2026}

The class denominator can be sharpened from the safe factorial 6 to 2. The character (SR6.2) has \(\operatorname{ch}_1=2H_0\) and \(\operatorname{ch}_2=0\), so the integral Chern classes of the perfect complex satisfy \(c_1=2H_0\) and \(c_2=2H_0^2\). Newton's formula gives \(\operatorname{ch}_3=-(2/3)H_0^3+c_3/2\). Comparing with (SR3.3) yields

\[
\begin{gathered}t_0=\frac{c_3(E_0^{\prime})}{2}+\left[\frac{16}{3}d^2N^6(1-8d)\right]\theta_{\rm prod}^3,\\\delta=2,\qquad T_0:=-\int_{E_i^6}t_0^2=32768d^4N^{12}.\end{gathered}\tag{SR6.3}
\]

The bracket is an integer because 3 divides \(L\). Hence \(2t_0\) is integral. The same calculation expresses 2t as an integral algebraic cycle on the ordinary-deformation open. The degree bound of its positive and negative parts can be chosen uniformly on a nonempty finite-type open by representing the relative Chern classes and the polarization power by cycles there. The proper relative Chow-space argument of supplied DIV1 then extends this fixed multiplier 2, with some common finite degree bound, to every fiber of the fixed family. This argument establishes a uniform multiplier; it does not calculate the degree bound.

For completeness, the exact norm in (SR6.3) is obtained before the isogeny. The two decomposable determinant vectors square to zero. In the Darboux orientation of D-SR-03 their product has top coefficient \((s-\bar s)^6\). The product of the \(\theta\) and dual-theta symplectic volumes is one. Therefore

\[
\begin{gathered}\int_A t^2=2\Delta^2(a\bar a)^2(s-\bar s)^6=-8d,\\-\int_{E_i^6}(f^*t)^2=8d\deg f=32768d^4N^{12}.\end{gathered}\tag{SR6.4}
\]

Both the twelve-generator wedge calculation and the two-by-two B-field and polarization matrices are checked by \texttt{semiregular\_gaussian\_certificate.py}, with their exact outputs in its JSON file. These arithmetic checks do not claim that the large torsion subgroup was enumerated.

There is also a fully integral normalization for every field in D-SR-05 without diagonalizing its action over \(\mathbb Q(i)\). Use the same Fermat map \(h\) and rational k, put \(M=2dN\), and form \(f_m=Mk:E_i^6\to A\). It is an actual isogeny because \(M\) is divisible by \(N\). Let \(\Theta_S\) be the integral ample class on \(E_i^6\) with Hermitian matrix \(S=\bigl(\begin{smallmatrix}N&1\\1&1\end{smallmatrix}\bigr)\) tensor \(I_3\). Then

\[
\deg f_m=64M^{12},\qquad f_m^*H=2M^2\Theta_S,\qquad B_m=-\tfrac12f_m^*\gamma=-M^2C_2\otimes I_3.\tag{SR6.5}
\]

The matrix \(B_m\) is integral. The K-action on the source is the integral scalar-block matrix \(T+1\), and its Hermitian space is split by D-SR-05. Twist \(f_m^*E\) by the line bundle of class \(B_m\), and conjugate its deck action exactly as above. The negative-self-Ext proof of D-SR-04 applies to this isogeny as well. The full ordinary character is \(t_m+2H_m-H_m^3/(12d)+H_m^5/(960d^2)\), where \(t_m=f_m^*t\) and \(H_m=2M^2\Theta_S\). The same Chern-class computation gives

\[
\begin{gathered}t_m=\frac{c_3}{2}+\left[\frac{128}{3}d^5N^6(1-8d)\right]\Theta_S^3,\\\delta_m=2,\qquad T_m:=-\int_{E_i^6}t_m^2=512dM^{12}=2^{21}d^{13}N^{12}.\end{gathered}\tag{SR6.6}
\]

The bracket is an integer since 3 divides \(d\). This gives an ordinary CM seed and explicit class normalization for every imaginary quadratic \(K\). Its underlying elliptic factor still has CM by \(\mathbb Q(i)\), while its Weil action is \(K\). The balanced standard \(E_K^6\) assertion is made only for the Gaussian case of D-SR-04. The polarization \(\Theta_S\) and the product principal polarization in (SR6.3) are kept distinct.

\subsection{Quantitative scope and framework}

The independent review of D-SR-01--04 checked the full trace tensor-product argument, Fourier signs, rational volume, and ordinary negative self-Ext. Exact source locators are recorded in \texttt{SOURCES.md}: Buchweitz--Flenner2003 Proposition4.2 pp166--167; Buchweitz--Flenner2008 author-version Proposition6.2.1 and Theorem6.4.2 pp29--31; Mukai1981 Theorem2.2 p156 and equation3.1 p158; Beauville2014 Lemma2 p106 and Example3 p107; Perry2026v2 Theorem1.2 and Remarks2.3/2.5, Definition2.6/Lemma2.7. Poincare and very ampleness are classical inputs explicitly identified in that log. \cite{SR-Perry2026}

Bounding the cycle presentation after S-14 is a stricter question than existence of the semiregular seed. The present rank-zero proof answers the existence question; the separately checked SF-QC01--05 positive-rank construction and HD-CM6 lattice transfer answer the numerical question for the split class-number-one families. In the project's language, the object is now in Perf rather than only a numerical lambda-ring class, and the Hochschild trace verifies that its actual obstruction space is controlled by the Hodge character. Its nonzero weight-six determinant component is of type (3,3); the remaining components are polarization powers. The supplied bounded-division reduction turns its deformation into the integral uniform form of Grothendieck's condition A (cote16, scan050, corrected SIGN\_NEXT SN-C03/C04). The positive ample class \(H\) supplies compatibility with condition I (scan049 and cote44 Theorem4), but the cycles come from the categorical deformation theorem, not from positivity alone. The resultant/Cayley--Hamilton division step remains the supplied Theorem W, corresponding to scans048--049. No new historical attribution of perfect complexes or Hochschild semiregularity to those archive pages is intended.
